\section{Week 3: Purity of Will}

\begin{quotationx}
It is futile to attempt to be concentrated if the Will is passionate about other things. The oscillations of the mind
will never be able to achieve silence unless the Will itself infuses it with silence. Only the still Will can render
the imagination and the intellect silent in concentration.

\textbf{St. John of the Cross} and \textbf{St. Theresa d'Avila} never tire of repeating that the
concentration necessary for spiritual prayer is the fruit of the moral purification of the Will. 
\flright{\textsc{Valentin
Tomberg}, \emph{Meditations on the Tarot}}

\end{quotationx}

Concentration can be applied on three planes:

\begin{itemize}
\item Mental 
\item Astral 
\item Physical 
\end{itemize}
We began with learning concentration on the physical plane. Then we transferred that knowledge to our thoughts or mental
plane. Finally, we will do the same to our emotional life for the purification of the soul.

Note that there are many more levels beyond these. In the \emph{Letter on the Star}, Tomberg explains:

\begin{quotationx}
There are \emph{twelve} degrees higher than that of the consciousness of the human transcendental Self. It is necessary,
therefore, in order to attain to the ONE God, to elevate oneself successively to degrees of consciousness of the nine
spiritual hierarchies and the Holy Trinity. 

\end{quotationx}
The Mental, Astral, and Physical correspond to the spirit, soul, and body. In the Letter on Judgment, Tomberg relates
them to the Trinity. The undivided self, then, corresponds to the Unity of God.

\paragraph{Image and Likeness}
The idea of man being the “image and likeness” of God is a recurrent theme throughout Meditations on the Tarot. Although
people today often like to repeat that we are all born in the “image and likeness” of God, that is not at all the
Traditional teaching: rather, because of the Fall, we have lost the full likeness and it is the task of the Hermetist
to restore it. Tomberg explains:

\begin{quotationx}
The ideal of alchemical \emph{transformation} of Hermetism offers to human beings the way to the realisation of true
human nature, which is the image and likeness of God. Hermetism is the re-humanisation of all elements of human nature;
it is their return to their true essence. Just as all base metal can be transformed into silver and into gold, so are
all the forces of human nature susceptible to transformation into “silver” or “gold”, i.e. into what they \emph{are}
when they share in the image and likeness of God. 

\end{quotationx}
If we are already in the “image and likeness of God”, then our level of being as such right now is perfect: i.e., there
is no need for transformation, redemption, or regeneration.

The image of God, according to \textbf{St. Bernard}, is our “essential” being. In that case it must be our higher
intellectual soul, which distinguishes humans from animals. It is unsullied, it has no negative part, it is free, it is
the source of the “spark of God”, and so is perfect. However, we rarely live at that level of awareness. It is as
though we own a penthouse suite, yet choose to live in the basement.

The likeness, on the other hand, is our soul life which reflects the image. This is — because of
the various perturbations — what must be purified.

\paragraph{Emotions}
As was mentioned last time, personal emotions need to be silenced to make the soul capable of “receiving from above the
revelation of the word, the life and the light.”

Now, the emotional center of our being, or the “astral” plane, has its own way of knowing. This is called the “cognitive
power of the emotions”. This manner of knowing is quite different from that of the thinking center or mental plane.
This knowing is episteme, the knowledge of the heart, beyond the dianoia of mental knowing. There is a higher emotional
component concomitant with its knowing.

Our age is dominated by thinking, arguing, and so on. This dualistic thinking distorts the emotional center. Tomberg
writes this about the relationship between the will and thinking:

\begin{quotationx}
Thus, it is not thought as such which allows the desire for personal greatness or the tendency towards megalomania, but
rather \emph{the will }which makes use of the head and which can take hold of thought and reduce it to the role of its
instrument.

\emph{Organic} humility, replacing the current of the will-to-greatness is not found in the head, but rather in the
heart, i.e., it reaches the heart, penetrating from the right-hand side. Because it is there that the will-to-greatness
has its origin and it is there from whence it takes hold of the head and makes it its instrument. This is why many
thinkers and scientists want to think “without the heart” in order to be objective, which is an illusion, because one
can in no way think without the heart, the heart being the activating principle of thought; what one can do is to think
with a humble and warm heart instead of with a pretentious and cold heart. 

\end{quotationx}
When functioning well, the heart and the head cooperate. In the example of megalomania, on the other hand, we see that
the will can take hold of the head, making it the servant of a disordered emotion. Common knowledge warns us about
making decisions when in a negative emotional state, but that is often ignored. Moreover, it is even celebrated, since
an opinion stated with strong negative emotions is falsely given a higher value.

The other distortion is when the head tries to think without the heart under the guise of objectivity. This leaves our
emotional range limited and underdeveloped.

\paragraph{Purity of Heart is to Will One Thing}
The inner life of the soul, in our present condition, does not present a unity. Rather, our desires, aspirations,
passions, and so on, are in conflict with each other. First one dominates, then another, as though there were multiple
separate “I's” inhabiting, and even fighting for control over, the soul. Tomberg calls these “lost
sheep” alluding to the Gospel story. He explains:

\begin{quotationx}
The soul's faults and vices are not, fundamentally, monsters but rather, lost sheep. … As it is the
same with all the soul's faults and vices, we all have the mission of finding and bringing back to
the flock (i.e. to the soul's choral harmony) the lost sheep in ourselves. We are missionaries in
the subjective domain of our own soul, charged with the task of the conversion of our desires, ambitions, etc. We have
to \emph{persuade }them that they are seeking the realisation of their dreams in a false way by showing them the true
way. It is not a matter of commandment, but rather of the alchemy of the cross, i.e. making present an alternative way
for our desires, ambitions, passions, etc. It is a matter, moreover, of the alchemical “marriage of opposites”. 

\end{quotationx}
Our alchemical task, then, is the transmutation of these multiple selves into a single I.

\flright{\itshape Posted on 2022-12-09 by Cologero}
